\documentclass [a4paper, 10pt]{scrartcl}

% Sprache + Encoding
\usepackage{ucs} % unicode
\usepackage[utf8x]{inputenc}
\usepackage{ngerman} % deutsche Sprache
\usepackage{ragged2e}

\parindent0pt
\parskip5mm plus1mm minus1mm

\setlength{\textwidth}{140mm}
\setlength{\textheight}{265mm}
\setlength{\oddsidemargin}{12mm}
\setlength{\evensidemargin}{12mm}
\setlength{\topmargin}{-20mm}
\setlength{\headheight}{0mm}
\setlength{\columnsep}{5mm}
\setlength{\footskip}{18mm}
\setlength{\headsep}{12mm}

\usepackage{german}
\usepackage{graphicx}
\usepackage{tabularx}
\usepackage{amsmath}
\usepackage{mdwlist}
\usepackage{tabularx}



\newcommand{\Referent}[2][]{
  \newcommand{\betreuertitel}{#1\ }
  \newcommand{\betreuername}{#2}        
}


\newcommand{\Mitarbeiter}[2][]{%
  \newcommand{\mitarbeitertitel}{#1\ }%
  \newcommand{\mitarbeitername}{#2}%
}

\newcommand{\Student}[2][]{%
  \newcommand{\studenttitel}{#1\ }%
  \newcommand{\studentname}{#2}%
}
\def\Matrikelnummer #1 {%
  \def\matrikelnr {#1}%
}

\def\Titel #1 {%
  \def\arbeittitel {#1}%
}

\def\Arbeitsart #1 {% 
  \def\art {#1}
}

\newenvironment{aufgaben}[0]{\textbf{Aufgaben:}\\\vspace*{-10mm}\begin{itemize*}}{\end{itemize*}}

\def\Termine #1#2#3 {%
  \def\beginn {#1} \def\zwischen {#2} \def\ende {#3}%
}

\def\Anfang{%
  \centerline{\resizebox{\linewidth}{!}{\includegraphics{isasheader.pdf}}}%
  % 
  \sffamily 
  \vspace*{5mm}%
  \centerline{\Large{%
      \textbf{\art:}}}\vspace*{-5mm}
  \begin{center}
    \Large{%
      \textbf{\arbeittitel}}
  \end{center}
  \bigskip
}			

\def\Aufgaben #1{%
  \def\aufgabenbesch {#1}
}

\def\Ende{%
  \vfill
  \hspace{-\parindent}%
  \begin{tabularx}{\textwidth}{lXlX}
    % \bf Titel: &\bf Effiziente Prädiktion nichtlinearer Systeme basierend auf der Approximation von Transitionsdichten mittels Gaussian Mixtures\\
    \textbf{Bearbeiter:} &\textbf{\studenttitel\studentname} &
    \textbf{Matrikelnummer: \matrikelnr}
    \\\ \\
    \textbf{Betr. Mitarbeiter:}&\multicolumn{2}{l}{\textbf{\mitarbeitertitel\mitarbeitername}}\\
    \textbf{Referent:}& \multicolumn{2}{l}{\textbf{\betreuertitel\betreuername}}\\
    \ \\\end{tabularx}
  \begin{tabular}{lr}
    \textbf{Beginn:} &\textbf{\beginn}\\
    \textbf{Zwischenvortrag:} & \textbf{\zwischen}\\
    \textbf{Abgabe:} & \textbf{\ende}
  \end{tabular}
  \par
  \vfill


  Karlsruhe, den \today\par
  \vspace{1.5cm}
  \noindent\begin{tabularx}{1.0\linewidth}[h]{cXcXc}
    \underline{\hspace{4cm}}&\ &\underline{\hspace{4.5cm}}
    &\ &\underline{\hspace{4cm}}\\
    \emph{\studentname} && \emph{\mitarbeitername} && \emph{\betreuername}
  \end{tabularx}  

}

\Arbeitsart{Masterarbeit}
\Titel{Kalmanfilterung mit Pseudo-Messungen aus Compressive Sensing von nicht-sparsen Signalen}
\Ziel{Reduzierung der Anforderung auf der Anzahl der Messungen zur Rekonstruktion nicht-sparser Signalen durch Kalman-Filtered Compressive-Sensing (KFCS)}

\Student[B.Sc.]{Zhao, Haibin}
\Matrikelnummer{2247796}
\Referent[Prof. Dr.-Ing.]{Uwe D. Hanebeck}
\Mitarbeiter[M.Sc.]{Funk, Christopher}
\MitarbeiterB[Dr.-Ing.]{Noack, Benjamin}

% Beginn Arbeit / Zwischenvortrag / Abgabe
\Termine{15. April 2020}{$\approx$15. Juli 2020}{15. Oktober 2020}

% Datum für Unterschrift:
%\renewcommand{\today}{20. Mai 2007}

% Ganttchart
\Ganttchart{gantt.png}

\begin{document}
\Anfang

Wenn ein Signal sparse ist, kann man Compressive-Sensing (CS) anwenden, um das Signal durch weniger Messungen rekonstruieren. Aber CS ist jedoch nicht für nicht-sparse Signale geeignet. Daher muss man ein neues Verfahren aufstellen, um CS anzuwenden. Mögliche Maßnahmen sind:
\begin{enumerate}
\item Betrachten die Änderung des Signals, anstatt das Signal selbst. Wenn die Änderung des Signals sparse ist, kann man das Verfahren zur Schätzung von Koeffizienten auf der Änderung von Koeffizienten anwenden. Im Vergleich zu der Anforderung von sparsen Signalen ist dies eine schwächere Anforderung, insbesondere bei kleinem zeitlichem Abtastintervall, kann es häufig erfüllt werden.
\item Kombinieren CS mit dem Systemmodell (Kalman-Filter). D.h. man kann das rekonstruierte Signal aus CS als Pseudo-Messungen in Kalman-Filter einsetzen. Und gleichzeitig werden Gewichtung dynamisch bestimmt: bei den realen Messungen werden niedrige Messunsicherheiten und bei den Pseudo-Messungen hohe Messunsicherheiten eingestellt.
\item Reduzieren die Dimension des zu rekonstruierenden Signals. CS wird nicht auf das ganze Signal, sondern auf alle möglichen Sensorstellen angewendet. Dadurch kann die Dimension deutlich verringert werden, wodurch die Sparsheit des Signals erhöht wird.
\item Sparse Coding Representation (SCR). Sucht man mithilfe des maschinellen Lernens die Basisfunktionen, auf denen das Signal sparse dargestellt werden kann.
\end{enumerate}

Da im obigen Prozess nur die Änderung des Signals geschätzt wird, führt dies zu einer kontinuierlichen Akkumulation von Fehlern. Daher sind einige Methoden erforderlich, um den Fehlern zu reduzieren, z.B.
\begin{enumerate}
\item Nach dem Kalman-Filter wird das geschätzte Signal in die Koeffizienten der Basisfunktionen transformiert, und wenn die Änderung der Koeffizienten zum nächsten Schritt berechnet werden, basiert es nicht auf dem geschätzten Koeffizienten aus dem vorherigen CS-Verfahren, sondern auf dem transformierten Koeffizienten aus dem Kalman-Filter.
\item Wenn die Änderung der Koeffizienten zum nächsten Schritt berechnet werden, basiert es nicht auf dem transformierten Koeffizienten aus dem Kalman-Filter, sondern auf der Prädiktion aus dem Systemmodell.
\end{enumerate}
Unter diesen ist das erste Verfahren besser für unbekannte Eingangsgröße geeignet, und das zweite ist besser für bekannte Eingangsgröße geeignet.

Darüber hinaus kann man bei jedem CS-Verfahren iterativ durchführen. D.h. man berechnet die Änderung der Koeffizienten iterativ, bis die Endbedingung erreicht. 
\\
\\
\\
\\
\\
\\

\begin{aufgaben}
\item Literaturüberblick über verwandte Ansätze zur Zustandsschätzung geben.
\item Erstellen von Testumgebung (Temperaturverteilung)
  \begin{itemize}
  \item Modalanalyse
  \item Finite Differenz Methode
  \end{itemize}
\item Grundlage für KF und CS
  \begin{itemize}
  \item Kalman-Filter
  \item Compressive Sensing
  \end{itemize}
\item Kalman-Filtered Compressive Sensing
  \begin{itemize}
  \item zufällige Messungen + Kalman-Filter
  \item sparse Änderung
  \item Dimension Reduktion
  \item Sparse Coding Representation
  \item Integration von KF und CS
  \item Iteratives Verfahren
  \end{itemize}
\item Schließen
  \begin{itemize}
  \item Präsentation
  \item Masterarbeit screiben
  \end{itemize}
\end{aufgaben}%
\Ende
\end{document}
